\documentclass[11pt,a4paper]{article}

%================================================
% ENCODAGE ET LANGUE
%================================================
\usepackage[utf8]{inputenc}
\usepackage[T1]{fontenc}
\usepackage[french]{babel}

%================================================
% MISE EN PAGE
%================================================
\usepackage{geometry}
\geometry{margin=1.55cm}

%================================================
% PACKAGE GTOLI
%================================================
\usepackage{../styles/gtoli}

%================================================
% INFORMATIONS DU DOCUMENT
%================================================
\title{\textbf{Statistiques}\\[0.2cm]
\large Titre du chapitre}

\author{GToli}
\date{\today}

%================================================
% DOCUMENT
%================================================
\begin{document}

\maketitle

%================================================
\section{Notion statistique concernée}
%================================================

\begin{cadreobjectif}
Dans ce chapitre, l’élève doit être capable de :
\begin{itemize}
    \item identifier les variables statistiques étudiées ;
    \item distinguer les types de variables ;
    \item choisir l’outil statistique adapté ;
    \item interpréter des résultats statistiques ;
    \item lire et construire des graphiques ;
    \item rédiger une conclusion statistique correcte.
\end{itemize}
\end{cadreobjectif}

%================================================
\section{Situation statistique de départ}
%================================================

\begin{cadrebleu}{Contexte statistique}
Cette section permet d’introduire :
\begin{itemize}
    \item la population étudiée ;
    \item l’échantillon ;
    \item les variables observées ;
    \item la problématique statistique.
\end{itemize}

\medskip

\textbf{Question statistique :}

\dotfill
\end{cadrebleu}

%================================================
\section{Vocabulaire statistique}
%================================================

\begin{cadretheorie}
\begin{itemize}
    \item \textbf{Population :} ensemble des individus étudiés.
    \item \textbf{Échantillon :} sous-ensemble observé.
    \item \textbf{Variable :} caractéristique mesurée.
    \item \textbf{Effectif :} nombre d’observations.
    \item \textbf{Fréquence :} proportion observée.
    \item \textbf{Moyenne :} indicateur de tendance centrale.
\end{itemize}
\end{cadretheorie}

%================================================
\section{Type de variable}
%================================================

\begin{cadregris}{Identification des variables}
\renewcommand{\arraystretch}{1.35}
\begin{tabularx}{\textwidth}{>{\bfseries}p{4cm}X X}
\toprule
Type de variable & Caractéristiques & Exemple \\
\midrule
Qualitative nominale & Catégories sans ordre & Couleur des yeux \\
Qualitative ordinale & Catégories ordonnées & Niveau de satisfaction \\
Quantitative discrète & Valeurs entières & Nombre d’enfants \\
Quantitative continue & Valeurs mesurables & Taille, masse \\
\bottomrule
\end{tabularx}
\end{cadregris}

%================================================
\section{Rappel théorique}
%================================================

\begin{cadretheorie}
Cette section contient :
\begin{itemize}
    \item les indicateurs statistiques ;
    \item les représentations graphiques ;
    \item les relations statistiques ;
    \item les formules ;
    \item les conditions d’utilisation ;
    \item les principes d’interprétation.
\end{itemize}

\medskip

\textbf{Point essentiel :} un résultat statistique doit toujours être interprété dans son contexte.
\end{cadretheorie}

%================================================
\section{Indicateurs statistiques}
%================================================

\begin{cadrevert}{Mesures importantes}
Exemples d’indicateurs statistiques :

\[
\bar{x}=\frac{\sum x_i}{n}
\]

\[
s^2=\frac{\sum (x_i-\bar{x})^2}{n}
\]

\[
f_i=\frac{n_i}{n}
\]

\medskip

\textbf{À préciser :}
\begin{itemize}
    \item signification des symboles ;
    \item unité éventuelle ;
    \item interprétation.
\end{itemize}
\end{cadrevert}

%================================================
\section{Choix du graphique}
%================================================

\begin{cadregris}{Quel graphique utiliser ?}
\renewcommand{\arraystretch}{1.35}
\begin{tabularx}{\textwidth}{>{\bfseries}p{4cm}X X}
\toprule
Situation & Variable & Graphique conseillé \\
\midrule
Répartition qualitative & Qualitative & Diagramme en barres \\
Pourcentages & Qualitative & Diagramme circulaire \\
Distribution continue & Quantitative continue & Histogramme \\
Lien entre deux variables & Deux quantitatives & Nuage de points \\
\bottomrule
\end{tabularx}
\end{cadregris}

%================================================
\section{Méthode statistique}
%================================================

\begin{cadremethode}
Pour traiter une situation statistique :

\begin{enumerate}
    \item identifier la population et les variables ;
    \item déterminer le type de variable ;
    \item choisir l’indicateur adapté ;
    \item organiser les données ;
    \item construire le graphique pertinent ;
    \item effectuer les calculs ;
    \item interpréter les résultats ;
    \item rédiger une conclusion statistique.
\end{enumerate}
\end{cadremethode}

%================================================
\section{Exemple guidé}
%================================================

\begin{cadreexemple}
\textbf{Situation :}

Insérer ici une situation statistique représentative du chapitre.

\medskip

\textbf{Analyse :}
\begin{itemize}
    \item Quelle est la population étudiée ?
    \item Quelle est la variable ?
    \item Quel type de variable observe-t-on ?
    \item Quel indicateur faut-il utiliser ?
    \item Quel graphique est adapté ?
\end{itemize}

\medskip

\textbf{Résolution :}

Présenter ici les calculs et l’interprétation.
\end{cadreexemple}

%================================================
\section{Lecture de graphique}
%================================================

\begin{cadreenonce}
\textbf{Document statistique :}

Insérer ici un graphique, un tableau de données ou une représentation statistique.

\medskip

\textbf{Question :} analyser et interpréter les données.
\end{cadreenonce}

\begin{cadreaide}
Pour analyser un graphique :
\begin{itemize}
    \item identifier les axes ;
    \item repérer les unités ;
    \item observer les tendances ;
    \item comparer les valeurs ;
    \item interpréter le phénomène observé.
\end{itemize}
\end{cadreaide}

%================================================
\section{Solution détaillée}
%================================================

\begin{cadresolution}
La solution doit comprendre :
\begin{enumerate}
    \item l’identification des variables ;
    \item le choix des outils statistiques ;
    \item les calculs ;
    \item la construction ou lecture du graphique ;
    \item l’interprétation ;
    \item la conclusion statistique.
\end{enumerate}
\end{cadresolution}

%================================================
\section{Erreurs fréquentes}
%================================================

\begin{cadreattention}
Les erreurs fréquentes en statistiques sont :
\begin{itemize}
    \item confondre effectif et fréquence ;
    \item utiliser un graphique inadapté ;
    \item interpréter sans contexte ;
    \item oublier les unités ;
    \item confondre moyenne et médiane ;
    \item tirer une conclusion abusive.
\end{itemize}
\end{cadreattention}

%================================================
\section{Exercices progressifs}
%================================================

\begin{cadreactivite}
\begin{enumerate}
    \item identifier le type de variable ;
    \item calculer une moyenne ;
    \item calculer une fréquence ;
    \item construire un graphique ;
    \item interpréter un tableau ;
    \item rédiger une conclusion statistique.
\end{enumerate}
\end{cadreactivite}

%================================================
\section{Synthèse}
%================================================

\begin{cadresynthese}
À retenir :

\begin{itemize}
    \item variable étudiée : \dotfill
    \item type de variable : \dotfill
    \item indicateur principal : \dotfill
    \item graphique adapté : \dotfill
    \item erreur à éviter : \dotfill
    \item conclusion statistique : \dotfill
\end{itemize}
\end{cadresynthese}

%================================================
\section{Auto-évaluation}
%================================================

\begin{cadregris}{Je vérifie ma maîtrise}
\begin{itemize}
    \item[$\square$] Je sais identifier une variable.
    \item[$\square$] Je connais les types de variables.
    \item[$\square$] Je sais choisir un graphique.
    \item[$\square$] Je sais effectuer les calculs.
    \item[$\square$] Je sais interpréter un résultat.
    \item[$\square$] Je sais rédiger une conclusion statistique.
\end{itemize}
\end{cadregris}

\end{document}